% SOURCE_AUTHORITY: sga5_fr_workpass.tex, lines 8183--8230 (LF-normalized unit).
% SOURCE_SHA256: 791F4EFFC5E02832D5D77ED03518C8156D6F07E4C8238B03545DB93D883FBB28
\section*{Exposição VI}

\begin{center}
\textbf{COHOMOLOGIA \(\ell\)-ÁDICA}\\[0.5em]
por J.P. JOUANOLOU
\end{center}

Em toda a exposição fixa-se um número primo \(\ell\).

Como de costume, denota-se por \(X_{\mathrm{\acute et}}\) o sítio étale de um pré-esquema \(X\).

\subsection*{1. Feixes \(\ell\)-ádicos construtíveis}

Fixa-se de uma vez por todas, no presente número, um pré-esquema localmente noetheriano \(X\).

\subsubsection*{1.1. Definição e generalidades}

\begin{statement}{Definição 1.1.1}
Chama-se feixe \(\ell\)-ádico construtível sobre \(X\) a um sistema projetivo \((\ell\mathbb Z)\)-ádico de feixes abelianos construtíveis sobre \(X_{\mathrm{\acute et}}\) (V 3.2.1).
\end{statement}

Para que um sistema projetivo \((\ell\mathbb Z)\)-ádico
\[
F=(F_n)_{n\in\mathbb N}
\]
de feixes abelianos sobre \(X_{\mathrm{\acute et}}\) seja um feixe \(\ell\)-ádico construtível, basta que \(F_0\) seja construtível. Com efeito, a afirmação é local e, no caso noetheriano, segue-se de (V 5.1.2).

Denotemos por \(\Abc(X)\) a categoria dos feixes abelianos construtíveis sobre \(X\). Denotar-se-á por
\[
\ell\text{-}\adc(X)
\]
a subcategoria plena de \(\Hom(\mathbb N^{\circ},\Abc(X))\) engendrada pelos feixes \(\ell\)-ádicos construtíveis. Com a notação da Exposição V:
\[
\ell\text{-}\adc(X)=(\ell\mathbb Z)\text{-}\ad(\Abc(X)).
\]
Denotando por \(\Ab(X)\) a categoria dos feixes abelianos sobre \(X_{\mathrm{\acute et}}\), sabe-se (SGA 4 IX 2.9) que, quando \(X\) é noetheriano, os feixes abelianos construtíveis sobre \(X\) são os feixes abelianos noetherianos. Neste caso, a categoria \(\ell\text{-}\adc(X)\) identifica-se claramente com a categoria denotada por
\[
(\ell\mathbb Z)\text{-}\adn(\Ab(X))
\]
na Exposição V.

\begin{statement}{Proposição 1.1.2}
Seja \(F=(F_n)_{n\in\mathbb N}\) um objeto da categoria \(\Hom(\mathbb N^\circ,\Ab(X))\). Para que \(F\) seja \(\ell\)-ádico construtível, é necessário e suficiente que o seja localmente para a topologia étale de \(X\).
\end{statement}

\begin{prooftext}[Demonstração]
É claro que a propriedade de \(F\) de ser \(\ell\)-ádico se verifica sobre um recobrimento do objeto final do sítio. Por outro lado, de (SGA 4 IX 2.8) segue-se que os \(F_n\) são construtíveis se o são localmente.
\end{prooftext}
